\documentclass[11pt]{article}

\usepackage{nmrdoc}
\usepackage{siunitx}
\usepackage{bm}

\sisetup{per-mode=symbol}
\DeclareSIUnit{\angstrom}{\text{\AA}}
\DeclareSIUnit{\ppm}{ppm}

\begin{document}
\NmrTitle{The Larsen Empirical Hydrogen-Bond Shielding Calculator}

\section{What the calculator computes}

The Larsen hydrogen-bond calculator evaluates the empirical ProCS15
hydrogen-bond term for each atom in a protein conformation. It does not run an
electronic-structure calculation. It finds geometrically plausible
hydrogen-bond pairs, looks up precomputed DFT response grids indexed by that
geometry, rotates the returned \(3\times 3\) tensors into the protein frame,
and accumulates the atom-specific terms that Larsen's table assigns to each
backbone atom type.

The physical effect being represented is the change in nuclear shielding caused
when a donor hydrogen approaches an acceptor oxygen. A hydrogen bond polarizes
the local electronic structure of the donor and acceptor groups; those
electrons then change the secondary magnetic field at nearby nuclei. The
observable connected to this calculation is the shielding tensor
\(\bm\sigma\), usually reported experimentally after conversion to chemical
shift. In this implementation the grid tensors and the water offset are in
\(\si{\ppm}\), because the grid values are DFT shielding differences supplied
by the Larsen parameterization. The calculator maps structure to tabulated
tensors; it does not run the electronic response calculation or convert those
tensors into a complete chemical shift.

\section{Where shielding comes from}

For a nucleus placed in an applied magnetic field \(\bm B_0\), the local field
is written
\begin{equation}
  \bm B_{\mathrm{nuc}} = (\bm I-\bm\sigma)\bm B_0 .
  \label{eq:shielding}
\end{equation}
\(\bm I\) is the \(3\times 3\) identity matrix. The shielding tensor
\(\bm\sigma\) is dimensionless in physical NMR and is normally quoted in
parts per million. It is a linear-response coefficient: if the applied field is
changed a little, \(\bm\sigma\) tells how much of the induced electronic field
opposes or redirects that change at the nucleus.

The tensor form is needed because molecules are oriented objects. A field
applied along the peptide plane, for example, can induce a different secondary
field than a field applied normal to that plane. A matrix records this
directional response: the entry \(\sigma_{ab}\) says how an applied-field
component along Cartesian direction \(b\) contributes to the screened field
component along direction \(a\).

At the quantum-mechanical level, Eq.~\eqref{eq:shielding} comes from the
electrons responding to \(\bm B_0\). The perturbation changes their current and
spin response, and those electrons produce an induced magnetic field at the
nucleus. The Larsen calculator keeps a local empirical piece of that response:
the portion associated with hydrogen-bond geometry. Its input coordinates are
not electronic wave functions, but the hydrogen-bond distance and angles that
select a precomputed DFT difference tensor.

Solution NMR often reports an isotropic chemical shift, related to the average
shielding after reference subtraction and sign convention. The isotropic
number is not the whole tensor. Hydrogen bonding also changes the
orientation-dependent part of the shielding, and the code preserves that
rank-2 information rather than reducing the grid lookup to a scalar shift.

\section{The tensor split used by the code}

Every accumulated Larsen tensor is decomposed by the C++ type
\texttt{SphericalTensor}. For a real \(3\times 3\) tensor \(\bm X\), the scalar
part is
\begin{equation}
  T_0 = \frac{1}{3}\operatorname{tr}\bm X ,
  \label{eq:t0}
\end{equation}
the average of the three diagonal entries. This is the isotropic part; it is
the component that survives rotational tumbling.

The antisymmetric part is stored as a three-component pseudovector,
\begin{equation}
  \bm T_1 =
  \frac{1}{2}
  \begin{pmatrix}
    X_{yz}-X_{zy}\\
    X_{zx}-X_{xz}\\
    X_{xy}-X_{yx}
  \end{pmatrix}.
  \label{eq:t1}
\end{equation}
It stores the part of the matrix that changes sign under transpose. In routine
isotropic solution NMR this term is not usually observed directly, but the
code carries it because the grid returns full matrices.

The remaining symmetric traceless matrix is
\begin{equation}
  \bm S =
  \frac{\bm X+\bm X^{\mathsf T}}{2}-T_0\bm I,
  \qquad \operatorname{tr}\bm S=0 .
  \label{eq:s}
\end{equation}
It has five independent entries. The code stores those five entries in the
real spherical basis
\begin{equation}
  \bm T_2 =
  \begin{pmatrix}
    \sqrt{2}\,S_{xy}\\
    \sqrt{2}\,S_{yz}\\
    \sqrt{3/2}\,S_{zz}\\
    \sqrt{2}\,S_{xz}\\
    (S_{xx}-S_{yy})/\sqrt{2}
  \end{pmatrix}.
  \label{eq:t2}
\end{equation}
The factors in Eq.~\eqref{eq:t2} are the implementation's normalization. In
this basis, the squared length of the five-vector equals the Frobenius norm of
the matrix \(\bm S\):
\begin{equation}
  \sum_m T_{2m}^2=\sum_{ab}S_{ab}^2 .
  \label{eq:t2norm}
\end{equation}
The nine written doubles are ordered as \(T_0\), the three \(T_1\) components,
and the five \(T_2\) components. LarsenHBond populates all three pieces
whenever the grid tensor contains them. The separate water term is scalar in
\(\si{\ppm}\) and is written outside this tensor total.

\section{The method implemented}

The runtime first enumerates donor hydrogens. A backbone amide hydrogen
\(\mathrm{HN}(i)\) is an amide donor when the bond graph supplies the
preceding residue's carbonyl carbon; its donor-frame anchors are
\(\mathrm{N}(i)\) and \(\mathrm{C}'(i-1)\). An alpha hydrogen
\(\mathrm{H\alpha}(i)\), including both \(\mathrm{HA2}\) and
\(\mathrm{HA3}\) on glycine, is an alpha donor; its anchors are
\(\mathrm{C\alpha}(i)\) and \(\mathrm{N}(i)\). Candidate acceptor oxygen atoms
are found within \(\SI{4.2}{\angstrom}\) of the donor hydrogen, excluding
atoms in the donor's own residue.

An acceptor oxygen is then classified by its chemistry. Backbone carbonyl
acceptors use the NMA acceptor grid. Asn/Gln side-chain amide oxygens are
approximated by that same backbone-carbonyl grid, because the Larsen archive
has no separate side-chain-amide acceptor scan. Hydroxyl acceptors from Ser,
Thr, and Tyr use the methanol grid. Asp/Glu and terminal carboxylate oxygens
use the acetate grid, with the closer of the two carboxylate oxygens selected
as the accepting oxygen. These choices select one of six archives:
NMANMA, NMACOH, NMACOO, ALANMA, ALACOH, and ALACOO.

For a donor hydrogen position \(\bm h\), acceptor oxygen \(\bm o\), acceptor
heavy atom \(\bm c\), and acceptor-side third atom \(\bm t\), the lookup
coordinates are
\begin{equation}
  r = \|\bm o-\bm h\|,
  \qquad
  \theta =
  \cos^{-1}\!\left(
  \frac{(\bm h-\bm o)\cdot(\bm c-\bm o)}
       {\|\bm h-\bm o\|\,\|\bm c-\bm o\|}
  \right),
  \label{eq:rtheta}
\end{equation}
and
\begin{equation}
  \rho =
  \operatorname{atan2}\!\left(
    [\bm n_1\times\widehat{(\bm c-\bm o)}]\cdot\bm n_2,\,
    \bm n_1\cdot\bm n_2
  \right),
  \quad
  \bm n_1=(\bm o-\bm h)\times(\bm c-\bm o),\quad
  \bm n_2=(\bm c-\bm o)\times(\bm t-\bm c).
  \label{eq:rho}
\end{equation}
Here \(r\) is in \(\si{\angstrom}\), while \(\theta\) and \(\rho\) are in
degrees in the code. \(\theta\) is the angle H--O--C at the acceptor oxygen.
\(\rho\) is the signed H--O--C--third dihedral, using the IUPAC
\(\operatorname{atan2}\) convention. The Larsen archive filenames use the
opposite \(\rho\) sign, but the parser keyed the dense grids on the IUPAC sign;
the runtime therefore does not negate \(\rho\). The code accepts a geometric
hydrogen bond when \(\theta\in[\ang{90},\ang{180}]\). If \(\theta\) passes but
the lookup is outside the selected grid's stored \(r\) or \(\theta\) range, the
pair is counted as a grid miss; it still suppresses the water term for that
amide hydrogen because the geometry was hydrogen-bond-like.

The grid lookup is tensor-valued trilinear interpolation. For the selected
donor and acceptor classes, the archive stores tensors
\(\bm G_{ijk}^{(q)}\) on a regular grid in \(r,\theta,\rho\), where \(q\)
names the readout atom such as donor HN, donor C, or acceptor CA. The runtime
locates the enclosing cell and returns
\begin{equation}
  \bm G^{(q)}(r,\theta,\rho)
  =
  \sum_{\alpha,\beta,\gamma\in\{0,1\}}
  w_{\alpha\beta\gamma}\,
  \bm G^{(q)}_{i+\alpha,\,j+\beta,\,k+\gamma}.
  \label{eq:interp}
\end{equation}
The weights \(w_{\alpha\beta\gamma}\) are the usual products of the fractional
distances along the three axes and sum to one. The \(\rho\) axis wraps
periodically at \(\pm\ang{180}\); \(r\) and \(\theta\) must lie inside their
stored ranges, with a small floating-point tolerance at the bounds. The actual
metadata place the NMA-donor \(r\) axes near
\(\SIrange{1.5}{4.0}{\angstrom}\), the ALA-donor axes near
\(\SIrange{1.6}{4.0}{\angstrom}\) or
\(\SIrange{1.8}{4.0}{\angstrom}\) depending on acceptor class, and the
\(\theta\) axes near \(\ang{90}\) to \(\ang{180}\). Some dense-grid cells were
filled upstream for continuity; the calculator propagates a flag when any of
the eight interpolation corners was imputed.

The returned tensors are stored in a donor-centered frame. The frame rotation
is built from the donor atoms:
\begin{equation}
  \hat{\bm z}=\frac{\bm h-\bm a}{\|\bm h-\bm a\|},\qquad
  \hat{\bm x}=
  \frac{(\bm h-\bm u)-[(\bm h-\bm u)\cdot\hat{\bm z}]\hat{\bm z}}
       {\|(\bm h-\bm u)-[(\bm h-\bm u)\cdot\hat{\bm z}]\hat{\bm z}\|},
  \qquad
  \hat{\bm y}=\hat{\bm z}\times\hat{\bm x}.
  \label{eq:frame}
\end{equation}
\(\bm a\) is the donor anchor and \(\bm u\) is the donor-side third atom
listed above. Eq.~\eqref{eq:frame} is a Gram--Schmidt projection: it makes an
orthonormal basis by keeping the H--anchor direction as the \(z\)-axis and
removing that \(z\)-component from the H--third direction before making the
\(x\)-axis. With \(R\) storing \(\hat{\bm x},\hat{\bm y},\hat{\bm z}\) as its
rows, the protein-frame tensor is
\begin{equation}
  \bm\sigma_{\mathrm{lab}} = R^{\mathsf T}\bm\sigma_{\mathrm{canon}}R .
  \label{eq:rotate}
\end{equation}
This is the standard rank-2 change of basis: both matrix indices rotate.

The per-atom dispatch then follows the Larsen contribution table encoded in
\texttt{LarsenContribDispatch}. Amide-H donors contribute primary
\(\mathrm{HB}\) terms to their own residue and secondary \(\mathrm{HB}\) terms
to the NMA acceptor side; alpha-H donors contribute the corresponding
\(\mathrm{H\alpha B}\) terms. The implemented mask is: N receives
\(1^\circ\mathrm{HB}\), \(2^\circ\mathrm{HB}\), and
\(2^\circ\mathrm{H\alpha B}\); \(\mathrm{C\alpha}\) and \(\mathrm{C}'\) receive
\(2^\circ\mathrm{HB}\); \(\mathrm{C\beta}\) is excluded from the physical sum;
\(\mathrm{H\alpha}\) and HN receive all four pair terms. Backbone NMA acceptor
readouts route N, HN, CA, and HA to residue \(j+1\), while the acceptor C
readout routes to the accepting residue \(j\). Hydroxyl and carboxylate grids
do not carry acceptor-side readouts. Side-chain carbonyls reuse the NMA grid,
but the code has no backbone \(j+1\) residue to route to, so it does not apply
secondary terms for those acceptors either.

For atom \(a\), the tensor total written by the calculator is
\begin{equation}
  \bm\sigma_a^{\mathrm{pairs}} =
  \bm\sigma_a^{1^\circ\mathrm{HB}}+
  \bm\sigma_a^{2^\circ\mathrm{HB}}+
  \bm\sigma_a^{1^\circ\mathrm{H\alpha B}}+
  \bm\sigma_a^{2^\circ\mathrm{H\alpha B}} .
  \label{eq:sum}
\end{equation}
Each term is a sum of rotated grid tensors assigned to that atom. The
\(\mathrm{C\beta}\) readout is rotated and written as a diagnostic tensor,
but it is not included in Eq.~\eqref{eq:sum} and does not increment the
per-atom pair count. Amide hydrogens with no successful grid hit and no
post-\(\theta\)-gate grid miss receive a separate water offset
\begin{equation}
  \Delta\sigma_w = \SI{2.07}{\ppm}.
  \label{eq:water}
\end{equation}
The code stores this as a scalar field on HN atoms, not as an addition to
\(\bm\sigma_a^{\mathrm{pairs}}\).

There are no Larsen-specific numeric keys in
\texttt{data/calculator\_params.toml}. The TOML H-bond entries there
(\(\SI{50.0}{\angstrom}\) dipolar maximum distance,
\(\SI{3.5}{\angstrom}\) counting radius, and a two-residue sequence
exclusion) belong to other hydrogen-bond calculators. LarsenHBond instead
uses the compiled constants \(\SI{4.2}{\angstrom}\), \(\ang{90}\),
\(\ang{180}\), and \(\SI{2.07}{\ppm}\), plus the supplied grid axes.

\section{Inputs and output}

The runtime inputs are atom positions in \(\si{\angstrom}\), residue atom
slots, the bond graph, typed atom semantics, and the spatial atom index. It
does not consume MOPAC charges, AIMNet2 charges or embeddings, ff14SB charges,
APBS fields, or DSSP hydrogen-bond assignments. Those are separate external
data streams in the wider pipeline. LarsenResidue perception is also not in
this calculator's runtime path.

The external data used here are the six dense HDF5 grids under
\texttt{data/larsen\_hbond\_grids}, or the directory supplied through the
runtime \texttt{larsen\_hbond\_grids} path. They are externally supplied
tabulated Larsen parameterization data. This document treats them as lookup
tables and does not rederive how they were fit.

\texttt{WriteFeatures} emits eight files: total pair tensor, the four
per-class pair tensors, the \(\mathrm{C\beta}\) diagnostic tensor, the water
term, and the per-atom pair count. The six tensor files are \(N\times 9\)
tables in the \(T_0,T_1,T_2\) order of Eqs.~\eqref{eq:t0}--\eqref{eq:t2};
the water file is \(N\times 1\) in \(\si{\ppm}\); the count file is a
one-dimensional integer array.
The tensors are empirical, geometry-indexed shielding-difference tensors from
the Larsen grids, not a fresh prediction of the chemical shift. There is no
separate fitted scalar applied after lookup inside this calculator; later
comparison with DFT reference shieldings belongs to the analysis that uses
these tensor features.

\end{document}
